\documentclass[12pt]{article}
\usepackage{bbm}
\usepackage{fancyhdr}
\usepackage{titling}
\usepackage{enumitem}
\usepackage{amsmath}
\usepackage{dsfont}
\usepackage{amsfonts}
\usepackage{amssymb}
\usepackage{amsthm}
\usepackage{geometry}
\usepackage{tikz}
\usepackage{algorithm}
\usepackage{algpseudocode}
\geometry{verbose,tmargin=1in,bmargin=1in,lmargin=1in,rmargin=1in}
\renewcommand{\vec}[1]{\mathbf{#1}}
\newcommand{\dec}{\text{dec}}
\newcommand{\divi}{\mathbin{\textnormal{div}}}
\title{Homework 1}
\author{Satvik Prasad}
\date{}
\newtheorem{lemma}{Lemma}[section]
\fancypagestyle{plain}{%
  \fancyhf{}
  \fancyhead[L]{\theauthor}
  \fancyhead[R]{\thedate}
}
\pagestyle{plain}
\makeatletter
\def\@maketitle{%
  \newpage
  \null
  \vskip 1em%
  \begin{center}%
    \let \footnote \thanks
    {\LARGE \@title \par}%
    \vskip 1em%
  \end{center}%
  \par
\vskip 1em}
\makeatother
\usepackage{cmbright}

\newcommand{\mat}[1]{\mathbf{#1}}

% ---------- Probability theory ----------
% Core operators (optional argument-free; use as \E{X}, \Var{X}, \Cov{X}{Y})
\DeclareMathOperator{\EOp}{\mathbb{E}}
\DeclareMathOperator{\VarOp}{Var}
\DeclareMathOperator{\CovOp}{Cov}
\DeclareMathOperator{\CorrOp}{Corr}

\renewcommand{\P}[1]{\mathbb{P}\left(#1\right)}        % \P{A}
\newcommand{\E}[1]{\EOp\left[#1\right]}                % \E{X}
\newcommand{\Var}[1]{\VarOp\left(#1\right)}            % \Var{X}
\newcommand{\Cov}[2]{\CovOp\left(#1,\,#2\right)}       % \Cov{X}{Y}
\newcommand{\Corr}[2]{\CorrOp\left(#1,\,#2\right)}     % \Corr{X}{Y}
\newcommand{\SD}[1]{\operatorname{SD}\left(#1\right)}  % \SD{X}

% Subscripted versions, e.g. \Esub{\theta}{X}, \Psub{\theta}{A}
\newcommand{\Esub}[2]{\EOp_{#1}\left[#2\right]}
\newcommand{\Psub}[2]{\mathbb{P}_{#1}\left(#2\right)}

% Conditioning, indicators, independence
\newcommand{\given}{\,\middle|\,}                      % \E{X \given Y}
\newcommand{\ind}[1]{\mathbbm{1}\left\{#1\right\}}     % \ind{X > 0}
\newcommand{\indep}{\perp\!\!\!\perp}                  % X \indep Y
\newcommand{\iid}{\overset{\text{iid}}{\sim}}          % X_i \iid \Normal(0,1)

% Convergence
\newcommand{\convp}{\xrightarrow{\;p\;}}               % in probability
\newcommand{\convd}{\xrightarrow{\;d\;}}               % in distribution
\newcommand{\convas}{\xrightarrow{\text{a.s.}}}        % almost surely
\newcommand{\convLp}[1]{\xrightarrow{L^{#1}}}          % in L^p

% Probability spaces
\newcommand{\F}{\mathcal{F}}                           % sigma-algebra
\newcommand{\B}{\mathcal{B}}                           % Borel sigma-algebra
\newcommand{\pspace}{(\Omega, \F, \mathbb{P})}

% Common distributions
\newcommand{\Normal}{\mathcal{N}}
\newcommand{\Bern}{\operatorname{Bernoulli}}
\newcommand{\Bin}{\operatorname{Binomial}}
\newcommand{\Pois}{\operatorname{Poisson}}
\newcommand{\Unif}{\operatorname{Uniform}}
\newcommand{\Exp}{\operatorname{Exp}}
\newcommand{\Geom}{\operatorname{Geometric}}
\newcommand{\Gam}{\operatorname{Gamma}}
\newcommand{\Beta}{\operatorname{Beta}}

% Number sets
\newcommand{\R}{\mathbb{R}}
\newcommand{\N}{\mathbb{N}}
\newcommand{\Z}{\mathbb{Z}}


\begin{document}
\subsection{Writing Proofs}
\begin{enumerate}
  \item We prove that $\mathbb F^{\infty}$ is infinite dimensional by
    proving that any basis of $\mathbb F^{\infty}$ cannot have a
    finite number of elements. Note that the $0$ vector of $\mathbb
    F^{\infty}$ is the sequence consisting of all zeroes (zero of
    $\mathbb F$), since by our definition of addition, $0 + a = a$
    for any sequence $a,$ and a sequence $0$ satisfying $(0)_i = 0
    \in \mathbb F.$

    For the sake of contradiction, assume we have some basis of
    $\mathbb F^{\infty}$ with length $n$, $b_1, b_2, \ldots, b_n \in
    \mathbb F^{\infty}.$

    Consider the $n + 1$ sequences $a_1, \ldots, a_{n + 1}$ defined by
    $$
    (a_{i})_j =
    \begin{cases}
      0 & i \neq j,  \\
      1  & i = j
    \end{cases}
    $$
    for $1 \leq i \leq n + 1,$ $j \geq 1.$

    Assume for the sake of contradiction that this set of vectors is
    linearly dependent. Therefore, we can conclude there exist
    $\mu_1, \ldots, \mu_{n+1} \in \mathbb F$ such that
    $$
    \mu_1 a_1 + \mu_2 a_2 + \ldots + \mu_{n+1} a_{n+1} = (\mu_1,
    \mu_2, \ldots, \mu_{n + 1}, 0, 0 \ldots) = \vec 0,
    $$
    and $(\mu_1, \ldots, \mu_{n + 1}) \neq 0 \in \mathbb F^{n + 1}.$

    Since $\mathbb F^{\infty}$ is a vector space, it has a unique
    zero vector. Thus, $(\mu_1, \ldots, \mu_{n + 1}, 0, \ldots)$ is
    composed of only $0$ entries, meaning $(\mu_1, \ldots, \mu_{n +
    1}) = 0 \in \mathbb F^{n + 1},$ which is a contradiction. Hence,
    $a_1, \ldots, a_{n + 1}$ are linearly independent.

    However, the length of a linearly independent list is at most the
    length of a spanning list. Since $b_1, \ldots, b_{n}$ are
    spanning, $a_1, \ldots, a_{n + 1}$ cannot be linearly
    independent, deriving a contradiction.

    Hence, we cannot have a basis of length $n$ for any integer $n,$
    meaning $F^{\infty}$ is infinite dimensional.

  \item Construct the map $T: \mathcal P_m(\mathbb F) \to \mathbb F$ defined by
    $$
    Tp = p(1).
    $$

    Note that $T(p + q) = (p + q)(1) = p(1) + q(1) = Tp + Tq,$
    $T(\lambda p) = (\lambda p)(1) = \lambda p(1) = \lambda T(p)$ and
    $Tc = c$ for $c \in \mathbb F,$ meaning $T$ is a surjective
    linear map, mapping from a space with dimension $m + 1$ to a
    space with dimension $1.$

    Hence, the null space of $T$ has dimension $m$ by the
    rank-nullity theorem. The null space consists of all $p \in
    \mathcal P_m(\mathbb F)$ such that $p(1) = 0.$ Thus, the $m + 1$
    vectors $p_0, \ldots, p_m$ reside in a space of dimension $m,$
    and thus they must be linearly dependent, as required.

  \item First, we prove that the vector space $V = \{\vec 0\},$ i.e
    consisting of only the zero vector, has a single basis. First,
    $\{\}$ is a basis for $V,$ since $\text{span}(\emptyset) = \{0\}
    = V,$ and $\emptyset$ is a linearly independent set.

    Assume $V$ has some other basis $B.$ Since $\vec 0$ is the only
    element in $V,$ this basis $B$ must contain $\vec 0$ at least
    once, i.e., $B_i = \vec 0.$ However, then a non-zero linear
    combination of the basis $B$ with the coefficients $\lambda_j =
    \delta_{ij}$ will produce the zero vector, meaning $B$ cannot be a basis.

    Hence, $V$'s only basis is $\emptyset,$ meaning $\{\vec 0\}$ is a
    vector space with a single basis.

    Then, assume we have any other vector space $U \neq \{\vec 0\}$
    over the field $\mathbb R$ with a single basis $B.$ Since
    $\text{span}(\emptyset) = \{\vec 0\} \neq U,$ we have $B \neq
    \emptyset.$ By definition,
    $$
    \text{span}(B) = U,
    $$
    and $B$ is linearly independent. Then, consider the set $B' =
    \{2b : b \in B\}.$ Note that $B \neq \emptyset,$ and $\vec 0 \not \in B.$

    Let $u \in U.$ Since $B$ spans $U,$ there exist $b_1, \ldots, b_k
    \in B$ and $\lambda_1, \ldots, \lambda_k \in \mathbb R$ such that
    $$
    u = \sum_{i=1}^{k} \lambda_i b_i = \sum_{i=1}^{k}
    \frac{\lambda_i}{2}(2b_i) \in \text{span}(B'),
    $$
    so $U \subseteq \text{span}(B').$ Additionally, since $B'
    \subseteq U$ and $U$ is closed under linear combinations,
    $\text{span}(B') \subseteq U,$ meaning $\text{span}(B') = U.$

    Additionally, because $B$ is linearly independent,
    $$
    \lambda_1 b_1 + \lambda_2 b_2 + \ldots = \vec 0
    $$
    has a single solution, which is the trivial solution. Assume that
    $B'$ is linearly dependent. Then,
    $$
    \lambda_1 (2b_1) + \lambda_2 (2b_2) + \ldots =   2(\lambda_1 b_1
    + \lambda_2 b_2 + \ldots) = \vec 0
    $$
    has a non-trivial solution, meaning
    $$
    \lambda_1 b_1 + \lambda_2 b_2 + \ldots = \vec 0
    $$
    has a non-trivial solution, which is a contradiction. Thus, $B'$
    is linearly independent.

    Therefore, we can conclude $B'$ is also a basis, meaning $B' =
    B.$ Therefore, for all $b_i,$ there exists some $b_j$ such that
    $$
    2b_i = b_j \Rightarrow 2b_i - b_j = 0.
    $$
    \begin{itemize}
      \item If $i = j,$ we must have $2b_i = b_i,$ which is a
        contradiction since $b_i \neq 0$ for all indices $i.$

      \item Otherwise, consider a single such $(i, j)$ pair where $i
        \neq j$. Then the non-trivial linear combination on $B$ with
        coefficients
        $$
        \lambda_i = 2, \lambda_j = -1, \lambda_{k\neq i, j} = 0
        $$
        produces the zero vector, which is a contradiction.
    \end{itemize}

    Thus, $B' \neq B,$ and there are $\geq 2$ basis sets on $U,$
    deriving a contradiction. Thus, for $U$ to have a single basis,
    it's basis set must be $\emptyset.$

  \item Suppose that $V$ is finite dimensional, and we have subspaces
    $U, W$ such that $U + W = V.$ Consider a basis of $U$, $B_U =
    \{u_1, \ldots, u_m\}$, and a basis of $W$, $B_W = \{w_1, \ldots, w_n\}.$

    Since we have $U + W = V$, for any $v \in V$, there exist $u \in
    U$ and $w \in W$ such that $u + w = v.$ Then, by the definition
    of a basis, there exist coefficients $\lambda_1, \ldots,
    \lambda_m$ and $\mu_1, \ldots, \mu_n$ such that
    $$
    \lambda_1 u_1 + \ldots + \lambda_m u_m + \mu_1 w_1 + \ldots + \mu_n w_n = v,
    $$
    and hence the set
    $$
    B' = \{u_1, \ldots, u_m, w_1, \ldots, w_n\}
    $$
    has $\operatorname{span}(B') = V.$

    Now, consider
    $$
    B = \{B'_i : B'_i \notin \operatorname{span}(B'_1, \ldots, B'_{i-1})\}.
    $$
    For linear independence, suppose
    $$
    \sum_{j=1}^{r} \lambda_j B'_{i_j} = 0
    $$
    for indices $i_1 < \cdots < i_r$ such that $B'_{i_j} \in B.$
    Consider the largest $k$ such that $\lambda_k \neq 0.$ Then, we can write
    $$
    \sum_{j=1}^{k} \lambda_j B'_{i_j} = 0
    \Rightarrow
    \lambda_k B'_{i_k} = -\sum_{j=1}^{k-1} \lambda_j B'_{i_j},
    $$
    and finally
    $$
    B'_{i_k} = -\frac{1}{\lambda_k}\sum_{j=1}^{k-1} \lambda_j B'_{i_j}.
    $$
    Letting $\mu_j = -\dfrac{\lambda_j}{\lambda_k}$, we have
    $$
    B'_{i_k} = \sum_{j=1}^{k-1} \mu_j B'_{i_j}
    \Rightarrow
    B'_{i_k} \in \operatorname{span}(B'_{i_1}, \ldots, B'_{i_{k-1}})
    \subseteq
    \operatorname{span}(B'_1, \ldots, B'_{i_k-1}),
    $$
    which is a contradiction of the definition of $B.$ Thus, $B$ is
    linearly independent.

    It remains to show that $\operatorname{span}(B) = V.$ We show
    that every $B'_i \in B'$ lies in $\operatorname{span}(B)$ by
    induction on $i.$ Suppose that
    $$
    B'_1, \ldots, B'_{i-1} \in \operatorname{span}(B).
    $$
    If $B'_i \in B$, then clearly $B'_i \in \operatorname{span}(B).$
    Otherwise, by the definition of $B$,
    $$
    B'_i \in \operatorname{span}(B'_1, \ldots, B'_{i-1}).
    $$
    By the induction hypothesis, each $B'_1, \ldots, B'_{i-1}$ lies
    in $\operatorname{span}(B)$, and therefore
    $$
    B'_i \in \operatorname{span}(B).
    $$
    Thus, every element of $B'$ lies in $\operatorname{span}(B)$, so
    $$
    \operatorname{span}(B') \subseteq \operatorname{span}(B).
    $$
    Since $B \subseteq B'$, we also have
    $$
    \operatorname{span}(B) \subseteq \operatorname{span}(B').
    $$
    Therefore,
    $$
    \operatorname{span}(B) = \operatorname{span}(B') = V.
    $$
    Hence, $B$ is a basis of $V.$ Finally, since
    $$
    B \subseteq B' \subseteq U \cup W,
    $$
    every element of $B$ lies in $U \cup W.$

  \item Counterexample: $V = \mathbb R^{4},$ $v_1 = e_1, v_2 = e_2,
    v_3 = e_3, v_4 = e_4,$ $U = \text{span}(e_1, e_2, e_3 + e_4).$
    Clearly $e_1, e_2 \in U.$

    We show $e_3 \notin U.$ Suppose $e_3 = a e_1 + b e_2 + c(e_3 +
    e_4)$ for some $a, b, c \in \mathbb R.$ Comparing the third
    coordinate gives $c = 1,$ while comparing the fourth coordinate
    gives $c = 0,$ a contradiction. Similarly, $e_4 \notin U.$

    Since $e_1, e_2, e_3, e_4$ are linearly independent, $e_1, e_2$
    cannot be a basis for $U,$ since $e_3 + e_4 \in U,$ and $e_3 +
    e_4 \in \text{span}(e_1, e_2)$ would require
    $$
    \lambda_1 e_1 + \lambda_2 e_2 - e_3 - e_4 = \vec 0
    $$
    for some $\lambda_1, \lambda_2 \in \mathbb R,$ which is a
    non-trivial linear combination of $e_1, e_2, e_3, e_4$ equal to
    $\vec 0,$ contradicting their linear independence. Hence $v_1,
    v_2$ is not a basis of $U,$ so the statement is false.
\end{enumerate}

\subsection{Evaluating Proofs}
This proof is logically correct, since $V \oplus U$ is a direct sum
only if for all $v \in V$ and $u \in U,$  $v + u = 0$ requires $v =
0$ and $u = 0.$

\subsection{Using Computers}
\begin{center}
  \includegraphics[width=0.9\linewidth]{Screenshot 2026-09-15 at 7.52.27 pm.png}
\end{center}

This statement is true. The transformation from basis $v_1, \ldots,
v_m$  to $w_1, \ldots, w_m$ is the matrix $P$ with $1$'s only in the
lower triangular entries, and $0$'s everywhere else. Since this has a
non-zero determinant and $v_1, \ldots, v_m$ are linearly independent,
it follows that $w_1, \ldots, w_m$ are also linearly independent.
Since we know the dimension of $V$ must be $m,$ then $w_1, \ldots,
w_m$ must also be a basis.

\end{document}
